\section{Introduction}
The problem of strain hardening in metals is one of the last frontiers of classical physics. A first-principles solution to this problem would involve a dynamical theory of dislocations---the line defects that carry plastic deformation---capable of predicting hardening coefficients present in phenomenological relations for the flow stress \cite{kocksPhysicsPhenomenologyStrain2003}. These hardening coefficients in turn seem to be strongly tied to the dynamic self-organization of the dislocations into various spatial patterns \cite{sauzayScalingLawsDislocation2011}. Continuum dislocation dynamics, which models the transport of various dislocation density measures representing  the local dislocation state of the crystal, seems to be a promising way forward toward a first-principles theory of dislocation patterning and, eventually, strain hardening. However, the so-called statistical storage problem—the tendency for any continuum dislocation density measure to underdetermine the local plastic deformation—has been a perennial impediment to the development of these theories, a statement as true today as it was when Kr\"{o}ner pointed it out in his final outlook on the field \cite{Kroner2001}. 

Upon beginning to set to the problem of continuum dislocation dynamics, one is immediately faced with a choice: which density measures to use? Different approaches to statistical storage motivate different answers to this question and naturally lead to substantially different theories of continuum dislocation motion.
This pernicious difference makes it difficult for researchers, especially the practitioners of the respective theories, to recognize the worthwhile and complementary contributions these theories make to the physics of plasticity or how the theories are related. The current work is intended as a bridge between two such frameworks of continuum dislocation dynamics, the line bundle  \cite{andersonDislocationCorrelationsContinuum2024} and the higher order  \cite{Hochrainer2015} formulations. Hence, it is worth describing the statistical storage problem in more detail and how these formulations address it differently.

Continuum dislocation fields, by their nature, treat a collection of dislocations in a representative volume around a single field point. Each of these dislocations has a line direction $\boldsymbol{l}$ and a Burgers vector $\boldsymbol{b}$. Taken together, these are related to the incompatibility of a dislocation which gives rise to long-range mechanical fields. The incompatibility (or Kr\"{o}ner-Nye) tensor field is simply the net incompatibility of all dislocations in the local collection:
$$
\boldsymbol{\alpha}(\boldsymbol{r}) = \sum \boldsymbol{l}\otimes \boldsymbol{b}.$$
Its relationship to mechanics led some to develop dynamical theories of the evolution of the incompatibility tensor \cite{muraContinuousDistributionMoving1963}, and indeed some are still pursuing such theories \cite{Acharya2006,Arora2020}. However, one of the difficulties identified by Kr\"{o}ner \cite{Kroner2001} is simply that in the summation over any collection of dislocations, cancellation can occur in either of these vector quantities. The incompatibility tensor thus does not describe all of the dislocation content of a collection, only its so-called geometrically necessary dislocation (GND) content, which, in Nye's sense \cite{Nye1953}, represents the net dislocation content needed to describe the mismatch of crystalline orientations across this volume. The dislocations are the dynamic objects, and so incompatibility has to supplement the motion of the GND content with a phenomenological term guessing the evolution of the unaccounted dislocations \cite{Acharya2006}. 

The fact that dislocations belong to well-defined crystallographic slip systems (combinations of Burgers vectors and slip planes to which the line direction is generally confined) suggests that there are two forms of cancellation in the incompatibility tensor that should be distinguished. If the dislocation collection is made up of enumerated ‘species’ of dislocations belonging to distinct slip systems, cancellation might occur when we a) obtain the average line direction on a single slip system (this subcollection shares a Burgers vector), or b) when summing over the various slip systems to obtain the net incompatibility. The latter is a delightfully geometric crystallography problem amounting to exploring the nullspace of a projection matrix and is given a definitive treatment in face-centered cubic (fcc) crystals by Arsenlis and Parks \cite{Arsenlis1999}. However, this issue can be simply sidestepped by considering a separate dislocation density field pertaining to each slip system \cite{El-Azab2018}. The former issue, the cancellation of line directions within a slip system population, still remains a nuanced problem to say the least.

The single-slip statistical storage problem essentially boils down to the question of when it is relevant. This problem can be traced back to two-dimensional theories of straight edge dislocations in the investigation of x-ray diffraction line profile analysis \cite{wilkensTheoreticalAspectsKinematical1970,Groma1998}. Biting the bullet of the single-slip problem involves considering dislocation subcollections of different direction with separate density fields; in two dimensions, this simply involves partitioning the dislocation population into positive and negative dislocations \cite{Groma1998,Groma1999,Zaiser2001,Groma2003}. The fact that such a toy model was already a rich playground for the investigation of collective dislocation phenomena \cite{ispanovityEvolutionCorrelationFunctions2008,dusanispanovityCriticalityRelaxationDislocation2011,dusanispanovityAvalanches2DDislocation2014,Ispanovity2017,Wu2018,wuCellStructureFormation2018,wuCyclicloadingMicrostructurepropertyRelations2019,Ispanovity2020} is evidence of the subtlety of the single-slip statistical storage problem. However, in three dimensions, the space of dislocation orientations is no longer trivial: the distribution over positive and negative orientation must be replaced with a distribution over a continuous orientation space (the unit circle for planar dislocations as in the case of fcc crystals) \cite{El-Azab2000a}. It was soon realized that this further complicates matters by also needing to include a density measure related to the curvature of dislocations of a given orientation in the collection \cite{Hochrainer2007}. This high-dimensional theory was subsequently simplified by considering only low-order moments of the orientation distribution \cite{Hochrainer2015}, and a commensurate analytical treatment of the energetics also followed shortly after \cite{Zaiser2015}. While this is a powerful theoretical framework, it is complex. This complexity leads to difficulties in closing the theory simply due to the difficulty of comparing the theory to discrete dislocation data or even underlying dislocation processes \cite{SandfeldPo2015,Monavari2016,Monavari2018,Sandfeld2019,Sudmanns2019,Sudmanns2020,Song2021,Hochrainer2022}. Due to this difficulty, the current trajectory of this theory is towards large length scale, phenomenological crystal plasticity especially in relationship to phase field plasticity theories \cite{Groma2015,Groma2021,zhangContinuumDislocationDynamics2025}.

One might question if all this theoretical machinery is necessary, or equivalently, whether single-slip statistical storage is always relevant. If it is not, the orientation distribution is single-valued at all points in space (all dislocations are parallel), analogously to a smooth bundle of lines. This possibility is primarily motivated by length scale considerations; if the density is resolved on a small scale,  often taken to be ~10-100 nm, there is an elastic alignment mechanism as well as an annihilation mechanism for oppositely signed dislocations. The assumption that the vector density accurately captures the line orientation of the local dislocation collection has several benefits. First, it admits a more straightforward dynamical analogy with discrete dislocations (as the streamlines of the density field \cite{Lin2021a}), e.g., for considerations of dislocation reactions \cite{Weger2019,Lin2020,Vivekanandan2021,vivekanandanDataDrivenApproach2023} and average driving forces \cite{Anderson2021,andersonDislocationCorrelationsContinuum2024}. Moreover, because of the detailed orientation description and its close tie to the incompatibility tensor, the mechanics of finite deformation—--which are necessary to describe the lattice rotations central to theories of strain hardening \cite{kocksPhysicsPhenomenologyStrain2003,sauzayScalingLawsDislocation2011}—--and their influence on the dislocation dynamics are more suitably described using the vector density theory \cite{Starkey2020,Hochrainer2020}. Due to the fine spatial resolution of the theory, the line bundle framework is sometimes maligned as a ``pseudo-continuum theory'' with no benefit surpassing discrete dislocation dynamics, which operates at a similar length scale \cite{Monavari2018,Hochrainer2022,Groma2021}. However, the stable computational complexity \cite{Lin2021} and formulation of finite deformation mechanics \cite{Starkey2020,starkeyTotalLagrangeImplementation2022} allow it to push to higher strain regimes than discrete dislocation dynamics. Moreover, the theory is a rich dynamical system in its own right, showing promise of studying the onset of dislocation patterns \cite{Xia2015}. 

So how should these approaches be reconciled? It is well-known that the vector density dynamics are recoverable as a special case of the high-dimensional dynamics \cite{Sedlacek2010,Hochrainer2010}. If the question is one of length scales, where is the transition between the vector density theory (which is obviously correct in the limit of finer scales) and the higher-order theory (which is obviously necessary in the limit of coarser scales)? Because both theories make claims regarding the orientation distribution, this question should, in principle, be answerable by studying the orientation statistics of dislocations. The present work sets out to accomplish this comparison by defining local orientation distribution functions in a spatial averaging framework that makes length scale explicit (section \ref{sec:2_localcollections}), situating points of comparison between the line bundle and higher-order theories in a theoretical analysis of orientation fluctuations (\ref{sec:3_reduced_descriptions}), and then evaluating this relationship between the line bundle and higher-order theories from statistical analysis of orientation fluctuations in discrete dislocation configurations (\ref{sec:4_calculation}-\ref{sec:5_results}). Implications of the present work for the field of continuum dislocation dynamics are discussed in section \ref{sec:6_disc}.
